\section{Constructing certificates by regridding}
\label{sec:regridding}

We construct a certificate without knowing the minimizer. At each stage we
build new partitions around the point obtained at the preceding stage and
use subtree slopes evaluated at that same point. The partitions need not
refine one another. The proof controls the sum of squared copy errors;
it does not require every reconstructed coordinate to lie within the
current mesh width of the minimizer.

Throughout this section, the models satisfy \eqref{eq:model-contract},
the bag gradients are $M$-Lipschitz in the Euclidean norm, and
\eqref{eq:model-growth} holds with $g>0$. Models are continuous and convex
on their closed boxes. An exact local optimization oracle returns both
the minimum value and an attaining point on each nonempty box
intersection. Bag values and gradients are also available exactly.
These assumptions specify an oracle model; the number of oracle calls
does not bound their internal cost. Fixed coordinates of the original
domain have been removed, as in the model section.

\subsection{Graded shell partitions}

For a point $c\in\R^d$ and radius $h>0$, put
$Q_\infty(c,h)=\{v:\norm{v-c}_\infty\le h\}$. Distances from a set are
infimum distances, and widths are the largest coordinate side length.

\begin{lemma}[Shell construction]\label{lem:regridding-shell}
Let $X'$ be a nondegenerate box in $\R^d$ whose sides are at most $s_0$,
let $c\in X'$, $h>0$, and let $\theta=2^{-\mu}$ with an integer
$\mu\ge1$. Set $H=\max\{0,\ceil{\log_2(s_0/h)}\}$. There is a finite
closed-box partition $\Pi(c;h,\theta)$ of $X'$ with disjoint interiors,
at most $(H+1)(4/\theta)^d$ members, and
\begin{equation}\label{eq:regridding-shell}
 \width(Y)\le\max\{h,\theta\dist_\infty(Y,c)\}
 \le h+\theta\dist_\infty(Y,c)
 \qquad (Y\in\Pi(c;h,\theta)).
\end{equation}
Each member retains a shell-level index between $0$ and $H$ and is a
clipped grid cube. The level-$0$ grid has at most two intervals in each
coordinate. Each other level has at most $4/\theta$ intervals in each
coordinate. In dimension zero use the single cell, of width zero.
\end{lemma}

\begin{proof}
Cut $Q_\infty(c,h)$ at the coordinate hyperplanes through $c$ into its
$2^d$ cubes of side $h$; these form level $0$. For
$\ell=1,\ldots,H$, divide $Q_\infty(c,2^\ell h)$ into a grid of side
$g_\ell=\theta2^{\ell-1}h$, anchored at its lower corner. There are
$4/\theta$ positions per coordinate. Because $1/\theta$ is an integer,
the boundary of $Q_\infty(c,2^{\ell-1}h)$ lies on grid hyperplanes.
Retain the cubes that are not contained in that inner cube. Their
interiors are outside the inner cube, so their distance from $c$ is at
least $2^{\ell-1}h$, and their width is
$g_\ell\le\theta\dist_\infty(Y,c)$.

The levels tile the outer cube with disjoint interiors. It contains
$X'$ because $c\in X'$ and $2^Hh\ge s_0$. Intersect all cubes with
$X'$ and discard intersections having empty interior in $\R^d$.
Clipping decreases widths and increases distances. It preserves
coverage: every point of $X'$ is a limit of interior points of $X'$
off the finitely many grid hyperplanes, and one retained closed
intersection contains a subsequence and hence its limit. Thus the
discarded pieces do not remove a boundary point from the union.
Each level contributes at most $(4/\theta)^d$ members, including the
central level, proving the claim.
\end{proof}

The fixed-coordinate reduction is necessary here. If $X'$ had a fixed
coordinate, every clipped cube would have empty ambient interior. One
can instead work in its affine hull; removing fixed coordinates gives
the same construction with a simpler dimension convention.

\subsection{The algorithm and its constants}

Write
\begin{equation}\label{eq:regridding-constants}
\begin{aligned}
 C_0&=k(k-1)p, & B_0&=\frac{MC_0}{2}+\frac A4,
 &\eta&=\frac g{20},\\
 C&=12B_0+\frac{6kC_0M^2}{\eta},
 & B&=\max\left\{p,\frac{2C}{g}\right\}.
\end{aligned}
\end{equation}
Choose a dyadic $\theta=2^{-\mu}\le1/2$ satisfying
\begin{equation}\label{eq:regridding-theta}
 12C_0\theta^2\le1,\qquad
 \sqrt{12}\,k\sqrt{C_0}M\theta\le\frac\eta4,\qquad
 12kB_0\theta^2\le\frac\eta4.
\end{equation}
A restriction with a zero coefficient imposes no condition. In
particular, these formulas allow $k=1$, $M=0$, and $A=0$.

Start from any $x^{(-1)}\in X$, with incumbent value
$\UB=F(x^{(-1)})$. At stage $j=0,1,\ldots$, set
$h_j=s_02^{-j}$ and $c=x^{(j-1)}$, and perform the following steps.
\begin{enumerate}
\item Build the fresh partitions
$\mathcal L_t=\Pi(c_{V_t};h_j,\theta)$ and, for $t\ne r$,
$\mathcal P_t=\Pi(c_{S_t};h_j,\theta)$.
\item Use the common slopes $\lambda_t(c)$ from
\eqref{eq:model-slope}. Compute the maximal intercepts and root bound
$\LB_j$ by the dynamic program of \cref{prop:model-dp}. Store the
minimizing local points and child-cell choices.
\item Backtrack a minimizing configuration and set
$x_i^{(j)}=z_i^{\topbag(i)}$. Update the incumbent if
$F(x^{(j)})<\UB$, retaining this feasible point and its value.
\item Stop if $\UB-\LB_j\le\eps$, returning the current certificate
and the retained incumbent.
\end{enumerate}
The next center is always the newly reconstructed point, even when an
earlier point remains the incumbent. All local subproblems have a
nonempty compact domain and continuous convex objective, so their
minimum is attained. Finite minima and backtracking therefore exist.
Ties can be resolved arbitrarily.

Certificate validity holds for every dyadic $\theta$, whether or not
\eqref{eq:regridding-theta} holds. Those conditions prove the
construction and size bounds, whereas the stopping test uses only
valid lower and feasible upper bounds.

\subsection{Aggregate copy drift}

Fix a stage and an arbitrary configuration. Let $x$ be its consistent
point, and define
\begin{equation}\label{eq:regridding-aggregate}
\begin{split}
 b_t&=\width(B_t),\qquad d_t=\width(D_t)\quad(t\ne r),\\
 Q&=\sum_t b_t^2+\sum_{t\ne r}d_t^2,\qquad
 E=\sum_t\norm{z^t-x_{V_t}}_2^2,\qquad R=\norm{x-c}_2.
\end{split}
\end{equation}
Here $d_t$ is a width, not a dimension.

\begin{lemma}[Occurrence-based drift]\label{lem:regridding-drift}
Every configuration satisfies $E\le C_0Q$.
\end{lemma}

\begin{proof}
For each variable $i$, its occurrence bags form a rooted subtree $T_i$
of size $m_i\le k$. Its root is $\topbag(i)$. On an edge of this
subtree, the child copy belongs to $D_t$ and the parent copy belongs
to $(B_{\pi(t)})_{S_t}$. These boxes meet. Joining the two copies
through an intersection point gives
\[
 \abs{z_i^t-z_i^{\pi(t)}}\le b_{\pi(t)}+d_t.
\]
The parent copy need not itself belong to $D_t$.

Make a matrix whose rows are the occurrences in $T_i$ and whose
columns are the bag and separator widths on $T_i$. The row for $t$
has a $1$ at the parent-bag width and child-cell width for every edge
on the path from $\topbag(i)$ to $t$, and zero elsewhere. It maps
these nonnegative widths to upper bounds on $\abs{z_i^t-x_i}$.
Each row has $2\operatorname{depth}_i(t)$ ones. Its squared operator
norm is at most its squared Frobenius norm, which is
\[
 2\sum_{t\in T_i}\operatorname{depth}_i(t)
 \le m_i(m_i-1)\le k(k-1).
\]
To see the depth bound, order the vertices with parents preceding
children; the depth of the vertex in position $a$ is at most $a-1$.
Sum the resulting coordinate inequalities. A bag-width square is
counted at most $\abs{V_t}$ times, and a separator-width square at
most $\abs{S_t}$ times. Both dimensions are at most $p$, so
\[
 E\le k(k-1)\left(\sum_t\abs{V_t}b_t^2+
                   \sum_{t\ne r}\abs{S_t}d_t^2\right)
 \le k(k-1)pQ.
\]
No tree depth or degree is used.
\end{proof}

\begin{lemma}[Grading absorption]\label{lem:regridding-grading}
If the first condition in \eqref{eq:regridding-theta} holds, then
\begin{equation}\label{eq:regridding-grading}
 Q\le12Nh_j^2+12k\theta^2R^2.
\end{equation}
\end{lemma}

\begin{proof}
Apply \eqref{eq:regridding-shell} to a bag leaf containing $z^t$:
\[
 b_t\le h_j+\theta\norm{x_{V_t}-c_{V_t}}_2
                    +\theta\norm{z^t-x_{V_t}}_2.
\]
The same inequality holds for $d_t$ with the vectors restricted to
$S_t$. A variable with $m_i$ occurrences belongs to exactly $m_i-1$
separators. Thus the sum of the bag and separator squared center
differences is at most $(2k-1)R^2$.

The copy error vanishes on $V_t\setminus S_t$: every such coordinate
has its top occurrence at $t$. Consequently its bag-plus-separator
sum of squares is exactly $2E$. There are at most $2N-1$ boxes.
Using $(u+v+w)^2\le3(u^2+v^2+w^2)$ gives
\[
 Q\le6Nh_j^2+3(2k-1)\theta^2R^2+6\theta^2E.
\]
By \cref{lem:regridding-drift}, the last term is at most
$6\theta^2C_0Q\le Q/2$. Absorb it and use $2k-1\le2k$.
When $C_0=0$, the last term is zero and the same bound holds directly.
\end{proof}

\subsection{The error relative to the current center}

\begin{lemma}[Uniform configuration error]\label{lem:regridding-error}
Under \eqref{eq:regridding-theta}, every configuration, with slopes
$\lambda_t(c)$, satisfies
\begin{equation}\label{eq:regridding-error}
 \abs{F(x)-\Phi}\le\eta\norm{x-c}_2^2+CNh_j^2.
\end{equation}
\end{lemma}

\begin{proof}
Let $\operatorname{err}_t=a_t(z^t)-m_{t,B_t}(z^t)$, so
$0\le\operatorname{err}_t\le Ab_t^2/4$. The telescoping identity
of \cref{lem:model-telescope}, applied at $c$, gives
\[
 \Phi=F(x)+\sum_tT_t-\sum_t\operatorname{err}_t,
 \qquad
 T_t=a_t(z^t)-a_t(x_{V_t})-
         \inner{\nabla a_t(c_{V_t})}{z^t-x_{V_t}}.
\]
The Lipschitz-gradient Taylor bound and Cauchy--Schwarz imply
\[
 \abs{T_t}\le
 M\norm{x_{V_t}-c_{V_t}}_2\norm{z^t-x_{V_t}}_2
 +\frac M2\norm{z^t-x_{V_t}}_2^2.
\]
Since each variable occurs in at most $k$ bags,
\[
\begin{split}
 \abs{F(x)-\Phi}
 &\le M\sqrt{k}\,R\sqrt E+\frac M2E+\frac A4Q\\
 &\le M\sqrt{kC_0}\,R\sqrt Q+B_0Q.
\end{split}
\]
Substitute \eqref{eq:regridding-grading} and use
$\sqrt{u+v}\le\sqrt u+\sqrt v$ to obtain
\[
\begin{split}
 \abs{F(x)-\Phi}
 &\le\sqrt{12kC_0}M\sqrt N h_jR+12B_0Nh_j^2\\
 &\quad+
 \left(\sqrt{12}\,k\sqrt{C_0}M\theta
               +12kB_0\theta^2\right)R^2.
\end{split}
\]
Young's inequality bounds the first term by
$\eta R^2/2+6kC_0M^2Nh_j^2/\eta$. The last two conditions in
\eqref{eq:regridding-theta} bound the coefficient in parentheses
by $\eta/2$. This proves the claim.
\end{proof}

The gradient at $c$ need not vanish. Its linear copy terms cancel
algebraically. This is why the estimate applies to boundary minimizers
without a stationarity hypothesis or differentiability of subtree
value functions.

\subsection{Contraction, termination, and box counts}

\begin{theorem}[Constructive regridding]\label{thm:regridding-main}
For every $\eps>0$, the algorithm with a grading ratio satisfying
\eqref{eq:regridding-theta} returns a valid certificate and feasible
incumbent with $\UB-\LB\le\eps$. With the constants in
\eqref{eq:regridding-constants}, put
\begin{equation}\label{eq:regridding-horizon}
 J=\max\left\{0,\ceil{\log_2\left(s_0\sqrt{gBN/\eps}\right)}\right\}.
\end{equation}
The algorithm stops by stage $J$, and each completed stage satisfies
\begin{equation}\label{eq:regridding-distance}
 \norm{x^{(j)}-x^*}_2^2\le BNh_j^2,
 \qquad 0\le\UB-\LB_j\le gBNh_j^2.
\end{equation}
Writing $m=4/\theta$, its final number of leaves plus cells is at
most $2Nm^p(J+1)$, and the total number created through stopping is
at most $Nm^p(J+1)(J+2)$.
\end{theorem}

\begin{proof}
A consistent configuration at $x^*$ exists: choose leaves and cells
containing its restrictions. It has value at most $f^*$ because
the models are lower bounds and all copy differences vanish.
By \cref{prop:model-dp}, a minimizing configuration has
$\Phi=\LB_j\le f^*$. Put $e_j=\norm{x^{(j)}-x^*}_2^2$.
Quadratic growth and \eqref{eq:regridding-error} give
\[
 ge_j\le F(x^{(j)})-f^*
 \le F(x^{(j)})-\LB_j
 \le\frac g{20}\norm{x^{(j)}-x^{(j-1)}}_2^2+CNh_j^2.
\]
The squared triangle inequality implies
\begin{equation}\label{eq:regridding-contraction}
 e_j\le\frac{e_{j-1}}9+\frac{10C}{9g}Nh_j^2.
\end{equation}
For $j\ge1$, the inductive bound $e_{j-1}\le BNh_{j-1}^2$
and $h_{j-1}=2h_j$ give
$e_j\le(4B/9+10C/(9g))Nh_j^2\le BNh_j^2$, since $B\ge2C/g$.
At stage zero, coverage of the variables gives $n\le pN$, and
$e_{-1}\le ns_0^2\le BNh_0^2$. Thus the base coefficient is at
most $B/9+10C/(9g)\le2B/3$, which also proves the claim.

For $j\ge1$, the squared center difference is at most
$2(e_j+e_{j-1})\le10BNh_j^2$. At stage zero it is at most
$4BNh_0^2$. The incumbent includes the current point, so
\[
 0\le\UB-\LB_j\le F(x^{(j)})-\LB_j
 \le(gB/2+C)Nh_j^2\le gBNh_j^2.
\]
No monotonicity of the lower bounds was used. The definition of $J$
ensures $gBNh_J^2\le\eps$, proving termination and the gap guarantee.

At stage $j$, the shell level count is $j+1$, because
$h_j=s_02^{-j}$. \Cref{lem:regridding-shell} gives at most
$2Nm^p(j+1)$ leaves and cells, including zero-dimensional cells.
Summing this bound over $j=0,\ldots,J$ proves the cumulative count.
Earlier partitions can be discarded after retaining the new center
and the incumbent.
\end{proof}

For the largest admissible dyadic ratio, with
$\kappa=M/g$ and $a=A/g$, direct solution of
\eqref{eq:regridding-theta} gives
\begin{equation}\label{eq:regridding-conditioning}
 \frac1\theta=O\left(1+\sqrt{C_0}
       +k\sqrt{C_0}\,\kappa
       +\sqrt{k(C_0\kappa+a)}\right).
\end{equation}
Also $C/g=6C_0\kappa+3a+120kC_0\kappa^2$.
These are numerical conditioning factors. Bounded width does not
imply bounded occurrence: in an edge-bag decomposition of a star, the
central variable occurs in all bags. The theorem removes a separate
degree parameter while retaining its stated dependence on $k$.

\subsection{Touching incidences and local work}

\begin{proposition}[Shell-incidence count]\label{prop:regridding-incidence}
Let $q=\max_{t\ne r}\abs{S_t}$, with $q=0$ when there are no
nonroot bags. At stage $j$, both the own-separator incidences and
the parent-leaf/child-cell incidences can be processed with at most
$N3^qm^p(j+1)^2$ incidences of each type. The exact local
optimization-call count through stage $J$ is at most
\begin{equation}\label{eq:regridding-calls}
 N3^qm^p\frac{(J+1)(J+2)(2J+3)}6.
\end{equation}
With bag evaluations and local optimizations counted as oracle
calls, coordinate processing and affine arithmetic require
$O(pN3^qm^p(J+1)^3)$ counted operations through stage $J$.
\end{proposition}

\begin{proof}
Consider a bag partition in dimension $d\le p$ and a separator
partition in dimension $q'\le q$, projected into that bag. For
one pair of levels, each partition is a subset of an unclipped
grid with at most $m$ positions per coordinate. Compare the
unclipped side lengths. If the bag grid is finer, an interval of
its side length meets at most three intervals of the separator
grid, including endpoint contacts. Each of at most $m^d$ bag
cubes therefore meets at most $3^{q'}$ separator cubes.

If the separator grid is finer, each separator cube meets at
most $3^{q'}$ projected bag-grid positions. Each such position
has at most $m^{d-q'}$ private-coordinate fibers, and there are
at most $m^{q'}$ separator cubes. This again gives at most
$3^{q'}m^d$ incidences. Clipping cannot create an intersection
between previously disjoint cubes. The central grid satisfies
the same position bound. Dimension-zero separators have one
cell and satisfy the count directly.

There are $(j+1)^2$ level pairs. Summing the own-separator
incidences over nonroot bags and including one minimum per root
leaf gives at most $N3^qm^p(j+1)^2$ local programs. Summing the
parent-leaf/child-cell incidences over the $N-1$ tree edges
gives at most $(N-1)3^qm^p(j+1)^2$ child incidences. These
determine the minimum compatible child intercept for each
parent leaf. Child choices are separate finite minima, so
there is no product over the children.

For completeness, incidence enumeration need not use a dense
all-pairs table. In the first case, enumerate the finer bag
grid and compute, in each separator coordinate, the at most
three coarse intervals that meet it. In the second case,
enumerate separator cubes, their at most three coarse
positions per coordinate, and the private bag fibers.
Filter the resulting pairs by the exact shell membership and
clipped-box tests. The enumeration costs $O(p)$ per candidate;
its candidate count has the same bound just proved. It can be
streamed while updating minima.

Slopes are obtained in $O(Np)$ arithmetic per stage by a
bottom-up accumulation of bag gradients. For each leaf,
the sum of child slopes is an affine coefficient vector;
the corresponding child-intercept sum can be accumulated
from the incidence minima. The total number of parent-leaf/
child combinations is at most $(N-1)m^p(j+1)$, using the
uniform leaf bound. Thus this additional work is covered by
the stated count. Finally,
$\sum_{j=0}^J(j+1)^2=(J+1)(J+2)(2J+3)/6$, proving
\eqref{eq:regridding-calls} and the cumulative work bound.
\end{proof}

The actual separator dimension $q$ can equal $p$ when adjacent
bags coincide. If no bag is contained in its parent, then $q\le p-1$.
We retain the supplied decomposition and assignment; merging bags
could change the bag-gradient and model-error constants.

Counting bag evaluations as oracle calls is material. The number
of factors or terms assigned to a bag is not bounded by $N$ and $p$.
Their representation, model-generation, evaluation, and numerical
certification costs must enter an implementation bound. Moreover,
the leaf-plus-cell measure does not include separate proof objects
for every local incidence or the bit lengths of numerical entries.
The following arithmetic section supplies such bounds for one
explicit representation.

\subsection{Searching without growth or error constants}

\begin{proposition}[Budgeted parameter search]\label{prop:regridding-search}
Suppose all requested oracle calls terminate. An algorithm that does
not use $g$, $M$, or $A$ to select its grading ratio can still return a
valid certificate and incumbent of gap at most $\eps$.
Fix any sufficient index $\mu_*$ and a bound $W\ge\max\{2,2^{\mu_*}\}$
on the counted work of its run through stopping. Then a simple
search uses $O(W\log(2+W))$ counted operations before returning.
Oracle-internal costs remain excluded.
\end{proposition}

\begin{proof}
For rounds $r=1,2,\ldots$, start a fresh run for each
$\mu=1,\ldots,r$ from the same initial point. Give each run
a budget of $2^r$ counted operations, covering box construction,
incidence enumeration, scalar comparisons and arithmetic,
bag-value and gradient evaluations, and local optimization
calls. Enforce the budget before each counted operation.
There is no separate stage cap. Stop at the first completed
stage that passes its own gap test.

For every trial, validity supplies $\LB\le f^*$ and the retained
feasible incumbent supplies $f^*\le\UB$. A successful test
therefore proves the required guarantee independently of
the trial's grading ratio. A sufficiently large $\mu_*$
satisfies \eqref{eq:regridding-theta}; by
\cref{thm:regridding-main} its run terminates in finite counted
work. At round $R=\ceil{\log_2 W}$, one has $R\ge\mu_*$ and
$2^R\ge W$, so that round permits the entire sufficient run.
The total trial budgets through round $R$ are
\[
 \sum_{r=1}^Rr2^r=(R-1)2^{R+1}+2\le2R2^R=O(W\log(2+W)).
\]
Every round makes finite progress because it contains finitely
many budgeted operations and each requested call terminates.
\end{proof}

One can take $W$ to be a constant multiple of the work bound in
\cref{prop:regridding-incidence} for the sufficient run; its factor
$m^p$ ensures $W\ge2^{\mu_*}$. The returned trial may have a different
grading ratio and stopping stage. Its gap guarantee is unchanged,
and its actual partitions satisfy the shell count for those returned
parameters. The sharper final-box bound in
\cref{thm:regridding-main} refers to a run with a fixed sufficient
ratio; parameter search does not identify its returned trial with
that run. The returned certificate has at most $O(W)$ boxes under
the work-budget convention, since their emission is charged in its
last trial budget. Total work and final output are distinct measures.
